% Generated draft: requires mathematical and source review.
\documentclass[11pt]{article}
\usepackage[margin=25mm]{geometry}
\usepackage{fontspec}
\setmainfont{Latin Modern Roman}
\usepackage{amsmath,amssymb,mathtools,unicode-math}
\setmathfont{Latin Modern Math}
\usepackage{enumitem,needspace,xurl,hyperref}
\setlength{\parindent}{0pt}
\setlength{\parskip}{5pt}
\setlength{\emergencystretch}{3em}
\newcommand{\N}{\mathbb{N}}
\newcommand{\Z}{\mathbb{Z}}
\newcommand{\Q}{\mathbb{Q}}
\newcommand{\R}{\mathbb{R}}
\newcommand{\C}{\mathbb{C}}
\newcommand{\F}{\mathbb{F}}
\newcommand{\A}{\mathbb{A}}
\newcommand{\PP}{\mathbb{P}}
\newcommand{\E}{\mathbb{E}}
\newcommand{\eps}{\varepsilon}
\newcommand{\dd}{\,\mathrm{d}}
\begin{document}
\section*{UnsolvedMath ID: 30006294}
\textbf{Source problem code:} catalogue entry
\subsection{Definitions and mathematical statement}
The supplied record concerns two-dimensional flow over topography whose spatial scale is strictly smaller than the domain scale; no scale separation is assumed. Writing these scales schematically as \(\ell_{\mathrm{top}}\) and \(L\), respectively, the stated restriction is
\[
0<\ell_{\mathrm{top}}<L,
\]
without an assumption that \(\ell_{\mathrm{top}}/L\) is asymptotically small. The record does not specify how these scales are measured. A \emph{condensate} means a large-scale circulation arising at the domain scale. The expressions \emph{topographically pinned vortex}, \emph{vortex strength}, and \emph{vortex location} are not given mathematical definitions in the supplied record.

\textbf{Research problem.} Extend the selective-decay variational theory for flow over topography to predict the strength and location of vortices pinned to that topography. The supplied context includes the question of extending variational principles designed for energy-conserving systems to forced, dissipative, out-of-equilibrium flows. No governing equation, forcing or dissipation law, spatial domain, boundary conditions, admissible data, variational functional, constraints, or prediction criterion is specified. Consequently, the record determines a research objective but does not determine a fully quantified mathematical assertion.

\textbf{Supplied context and qualifications.} The source excerpt reports a branch of variational solutions corresponding to a condensate. For weak initial energy, it predicts organization at the topographic scale without a condensate; above a threshold initial energy, it predicts the emergence of a condensate; at large initial energy, it reports condensates of the type observed in two-dimensional turbulence without topography. It reports very good quantitative agreement between the predicted condensate amplitude and pseudo-spectral simulations with small hyperviscosity. The supplied dated assessment further states that forced-dissipative topography theory reproduces condensates and identifies pinned vortices numerically, while prediction of their strength and location remains an open question. These are supplied contextual claims, not hypotheses or results independently verified in this draft. The excerpt is truncated, and the connection between its opening question and its subsequent discussion is incomplete.
\subsection{Short English statement}
Develop selective-decay theory to predict how strong topographically pinned vortices are and where they occur, in a setting where the topographic scale is smaller than the domain scale without requiring scale separation. The supplied context includes forced and dissipative flows, but the equations, vortex observables, and prediction criteria are missing, so a fully quantified statement requires further source clarification.
\subsection{Sources}
\begin{itemize}
\item Local UnsolvedMath registry, \texttt{data/unsolvedmath/problems.json}, permanent ID \(30006294\), title ``Predicting Topographically Pinned Vortices,'' exact source code \texttt{catalogue entry}. The supplied record labels the problem open and its statement \texttt{corrected\_verified}; these are registry labels, not independent mathematical verification. No individual proposer is supplied; the proposed year is 2025.
\item \emph{Mathematical Advances in Geophysical Fluid Dynamics} (2025), Oberwolfach Reports, pp.\ 1215--1270, DOI \texttt{10.4171/owr/2025/24}. Supplied source and literature-source label: ``original OWR report.'' \url{https://doi.org/10.4171/owr/2025/24}. The supplied background describes Oberwolfach Reports as documenting workshops and their open-problem discussions.
\item Supplied literature-source attribution: ``Gallet, two-dimensional turbulence above topography.'' \url{https://www.cambridge.org/core/journals/journal-of-fluid-mechanics/article/twodimensional-turbulence-above-topography-condensation-transition-and-selection-of-minimum-enstrophy-solutions/654AD74779542F674C5A51FE5F82F28}. The record does not supply a full bibliographic citation or identify Gallet as the proposer of this problem.
\item Literature assessment dated exactly \texttt{2026-08-22}: the supplied curation labels the problem open and states that forced-dissipative theory reproduces condensates and identifies pinned vortices numerically, but does not yet predict their strength and location. This is supplied dated curation derived from the OWR material, not an independently verified literature assessment; its time-sensitive status and cited claims require independent verification.
\item The supplied excerpt cites reference [2] for condensates in two-dimensional turbulence without topography, but does not supply the corresponding bibliographic entry.
\end{itemize}
\end{document}
